\documentclass{article}
\input{~/studies/preamble}

\title{Graph Neural Network}
\author{Hyoung Joon, Kim}
\date{\today}

\begin{document}

\maketitle
\tableofcontents
\newpage

\section{Introduction}
    \textbf{Graph neural networks} are used to capture the information under the graph structure.
    Its main usage are:
    \begin{itemize}
        \item Graph classification: Classify the graph itself; ex. Molecule property prediction;
        \item Link prediction: Predict if edge exists between nodes; ex. Recommender system;
        \item Node / Edge classification: Classfy the nodes or edges of the graph; ex. FDS
    \end{itemize}

\section{Main Ideas}
    \subsection{Basic Definitions}
        \begin{defn}{Graph}{graph}
            The \textbf{graph} $\mathcal G=(\mathcal V, \mathcal E)$ consists of set of \textbf{nodes} (or vertices) denoted by
            $\mathcal V$, and a set of \textbf{edges} (or links) denoted by $\mathcal E$. Each nodes are
            indexed with $n=1,2,\dots,N$ where $N$ is the total number of nodes. Edges from $n$ to $m$ is denoted by
            $(n,m)$. If two nodes are linked by an edge, they are called \textbf{neighbors}, and set of all neighbors
            of node $n$ is denoted by $\mathcal N(n)$. 
        \end{defn}

        We use \textbf{adjacency matrix} to denote the connections between nodes. If the edge exists
        between node $n$ and $m$, then the element of adjacency matrix $A_{nm}=1$, else $A_{nm}=0$.
        If the graph is undirected, the adjacency matrix is symmetric.
    
    \subsection{Permutation Invariance}
        The graph structure is invariant under the permutation operation. That is, even if we change the
        numbering of the graph, the network prediction should not change. Or, the network prediction may change,
        applying the same permutation to the prediction without permutation. That is, the network prediction is
        equivariant under permutation. Precise definition is given below;
        
        \begin{defn}{Permutation Invariance and Equivariance}{perInvEquivar}
            Let \textbf{P} be a permutation matrix. Let $\textbf{y}(\textbf{X}, \textbf{A})$ be a network prediction defined on the
            node data matrix \textbf{X} and adjacency matrix \textbf{A}. If the network satisfies following condition, the
            network is \textbf{invariant} under permutation;
            \[
            \mathbf{y}(\tilde{\mathbf{X}}, \tilde{\mathbf{A}})=\mathbf{y}(\mathbf{X}, \mathbf{A})   
            \]
            where $\tilde{\mathbf{X}}=\mathbf{PX}$ and $\tilde{\mathbf{A}}=\mathbf{PA}\mathbf{P}^{\top}$.
            If the network satisfies following condition, the network is \textbf{equivariant} under permutation;
            \[
            \mathbf{y}(\tilde{\mathbf{X}}, \tilde{\mathbf{A}})=\mathbf{P}\mathbf{y}(\mathbf{X}, \mathbf{A}).
            \]
        \end{defn}

    \subsection{Neural Message Passing}
        To use the permutation invariance and equivariance as the inductive bias of our network, it is not a good idea to
        use ordinary operations or flattening the matrix. Instead we adopt the method of CNN, such as parameter sharing.
        \subsubsection{Graph Convolutional Network}
            The graph network layer should satisfy following requirements:
                \begin{itemize}
                    \item Handle variable-length inputs: Since the number of nodes and edges are different along graphs;
                    \item Must be differentiable;
                    \item Scale well: Since graphs can be very large, like social network;
                \end{itemize}

            We use the embedding vector of nodes, $\textbf{h}_n$, and use the neighbors of each nodes the collect the information
            and make pre-activations while satisfying invariance.

            \begin{algorithm}[H]
                \DontPrintSemicolon
                \SetArgSty{textnormal}
                \caption{Simple message passing neural network}
                \SetKwInOut{Input}{Input}
                \SetKwInOut{Output}{Output}

                \Input{Undirected graph $\mathcal G=(\mathcal V, \mathcal E)$
                \newline Initial node embeddings $\{\textbf{h}_n^{(0)}\}$
                \newline $\text{Aggregate}(\cdot)$ function
                \newline Update($\cdot$, $\cdot$) function}
                \Output{Final node embeddings $\{\textbf{h}_n^{(L)}\}$}

                \BlankLine
                \hrule

                \For{$l\in \{0, \dots, L - 1\}$}{
                    $\textbf{z}_n^{(0)} \leftarrow \text{Aggregate}\left(\left\{ \textbf{h}_m^{(l)}:m \in \mathcal N(n) \right\}\right)$\;
                    $\textbf{h}_n^{(l+1)} \leftarrow \text{Update}\left(\textbf{h}_n^{(l)},\textbf{z}_n^{(l)}\right)$
                }
                \Return{$\{\textbf{h}_n^{(L)}\}$}
                
            \end{algorithm}

        The Aggregate($\cdot$) function collects the information of neighbor nodes and make
        a vector $\textbf{z}$. It must be invariant under the ordering of the embedding vector.
        The Update($\cdot$, $\cdot$) function then update the given node with the vector $\textbf{z}$
        and make a new node embedding vector.

        \subsubsection{Aggregate Operator}
            The possible forms for the Aggregate function are:
            \begin{itemize}
                \item Simple sum: $\displaystyle \text{Aggregate}\left(\left\{ \textbf{h}_m^{(l)}:m \in \mathcal N(n) \right\}\right)=\sum_{m \in \mathcal N(n)} \textbf{h}_m^{(l)}$
                \item Divided sum: $\displaystyle \text{Aggregate}\left(\left\{ \textbf{h}_m^{(l)}:m \in \mathcal N(n) \right\}\right)=\frac{1}{|\mathcal N(n)|}\sum_{m \in \mathcal N(n)} \textbf{h}_m^{(l)}$
                
                    or $\displaystyle \text{Aggregate}\left(\left\{ \textbf{h}_m^{(l)}:m \in \mathcal N(n) \right\}\right)=\frac{1}{\sqrt{|\mathcal N(n)||\mathcal M(m)|}}\sum_{m \in \mathcal N(n)} \textbf{h}_m^{(l)}$
                \item Fully connected layer: $\displaystyle \text{Aggregate}\left(\left\{ \textbf{h}_m^{(l)}:m \in \mathcal N(n) \right\}\right)=\text{MLP}_{\bm{\theta}}\left(\sum_{m \in \mathcal N(n)} \text{MLP}_{\bm{\phi}}(\textbf{h}_m^{(l)})\right)$
                with learnable parameter $\bm{\theta}$ and $\bm{\phi}$.
            \end{itemize}
            
        \subsubsection{Update Operator}
            The possible forms for the Update function are:
            \begin{itemize}
                \item $\text{Update}\left(\textbf{h}_n^{(l)},\textbf{z}_n^{(l)}\right)=f(\textbf{W}_{self}\textbf{h}_n^{(l)}+\textbf{W}_{neigh}\textbf{z}_n^{(l)}+\textbf{b})$
                , where $f$ is a nonlinear activation function and $\textbf{W}_{self}$, $\textbf{W}_{neigh}$ and $\textbf{b}$ are the learnable weights and biases. 
            \end{itemize}

    \subsection{Node Classification}
        We can classify the nodes by applying softmax function at the last layer of the graph neural network and define a loss function as the cross-entrophy loss across all nodes and all classes.
        We distinguish three types of nodes for training:
        \begin{enumerate}
            \item $\mathcal V_{train}$: Labeled, included in the message-passing operations and used to compute the loss function for training;
            \item $\mathcal V_{trans}$: Unlabeled, included in the message-passing operations but not used to compute the loss function for training; Graph without transductive nodes can be considered as supervised learning.
            \item $\mathcal V_{induct}$: Unlabeled, and does not participate in the training session, but participate in message-passing during inference session and labels are predicted; Graph without inductive nodes can be considered as semi-supervised learning.
        \end{enumerate}
    \subsection{Edge Classification}
        We can classify the edges by defining additional edge properties. For example, to predict if the edge exists or not, we can
        define the probability $p(n,m)$ for the presence of an edge between node $n$ and $m$ as $p(n,m)=\sigma(\textbf{h}_n^{\top}\textbf{h}_m)$.
    \subsection{Graph Classification}
        We can classify the graph itself, by training with a set of labeled graphs $\{\mathcal G_1, \dots,\mathcal G_N\}$. We may use
        the final node embeddings to yield predictions about graphs. The simplest form of such approach is to take the sum of final node embeddings and pass through some activation function;
        \[
        \textbf{y}=\textbf{f}\left(\sum_{n \in \mathcal V}\textbf{h}_n^{(L)}\right)
        \]

\section{Extensions of Graph Neural Networks}
    \subsection{Graph Attention Networks}
        We can apply attention mechanism to graph networks, too! We may define the Aggregate function as following:
        \[
        \textbf{z}_n^{(l)}=\text{Aggregate}\left( \left\{ \textbf{h}_m^{(l)}:m \in \mathcal N(n) \right\} \right)
        =\sum_{m \in \mathcal N(n)} A_{nm}\textbf{h}_m^{(l)}
        \]
        where $A_{nm}$ are the attention weights.
        The attention weights may be defined as
        \begin{itemize}
            \item $\displaystyle{A_{nm}=\frac{\exp{(\textbf{h}_n^{\top}\textbf{W}\textbf{h}_m)}}{\sum_{k \in \mathcal N(n)}\exp{(\textbf{h}_n^{\top}\textbf{W}\textbf{h}_k)}}}$
                where \textbf{W} is a learnable parameter.
            \item $\displaystyle{A_{nm}=\frac{\exp{(\text{MLP}(\textbf{h}_n, \textbf{h}_m))}}{\sum_{k \in \mathcal N(n)}\exp{(\text{MLP}(\textbf{h}_n, \textbf{h}_k))}}}$
                where MLP has a single continuous output and invariant to the ordering of the node. If the MLP is shared across all nodes in graph, this aggregation function is equivariant under node reordering.
        \end{itemize}
    \subsection{Edge and Graph Embeddings}
        We can introduce edge embedding vectors to manage the graph whose data is associated with edges. We can also consider the graph embeddings, which relates to the whole graph.
        This additional embeddings allow us to make richer set of graph features, and a general neural message-passing algorithm.

        \begin{algorithm}[H]
            \DontPrintSemicolon
            \SetArgSty{textnormal}
            \caption{General massage passing neural network}
            \SetKwInOut{Input}{Input}
            \SetKwInOut{Output}{Output}

            \Input{Undirected graph $\mathcal G=(\mathcal V, \mathcal E)$ \newline
                Initial node embeddings $\{\textbf{h}_n^{(0)}\}$,
                Initial edge embeddings $\{\textbf{e}_{nm}^{(0)}\}$,
                Initial graph embedding $\{\textbf{g}^{(0)}\}$ \newline
                $\text{Aggregate}_{node}(\cdot)$ function \newline
                $\text{Update}_{node}(\cdot,\cdot,\cdot)$ function, $\text{Update}_{edge}(\cdot,\cdot,\cdot,\cdot)$ function, $\text{Update}_{graph}(\cdot,\cdot,\cdot)$ function
            }
            \Output{Final node embeddings $\{\textbf{h}_n^{(L)}\}$\newline
            Final edge embeddings $\{\textbf{e}_{nm}^{(L)}\}$\newline
            Final graph embedding $\{\textbf{g}{(L)}\}$}

            \BlankLine
            \hrule

            \For{$l \in \{0, \dots, L - 1\}$}{
                $\textbf{e}_{nm}^{(l+1)} \leftarrow \text{Update}_{edge}(\textbf{e}_{nm}^{(l)}, \textbf{h}_n^{(l)}, \textbf{h}_m^{(l)}, \textbf{g}^{(l)})$\;
                $\textbf{z}_n^{(l+1)} \leftarrow \text{Aggregate}_{node}\left(\left\{ \textbf{e}_{nm}^{(l+1)}:m \in \mathcal N(n) \right\}\right)$\;
                $\textbf{h}_n^{(l+1)} \leftarrow \text{Update}_{node}\left( \textbf{h}_n^{(l)}, \textbf{z}_n^{(l+1)}, \textbf{g}^{(l)} \right)$\;
                $\textbf{g}^{(l+1)} \leftarrow \text{Update}_{graph}\left( \textbf{g}^{(l)}, \{ \textbf{h}_n^{(l+1)}, \{\textbf{e}_{nm}^{(l+1)}\} \} \right)$
            }

            \Return{$\{\textbf{h}_n^{(L)}\}$, $\{ \textbf{e}_{nm}^{(L)} \}$, $\textbf{g}^{(L)}$}
        
        \end{algorithm}

    \subsection{Geometric Deep Learning}
        Consider the task of predicting the properties of a molecule. The molecule can be represented as a list of
        spatial coordinates and types. If we use neural network for prediction, the network must be invariant of reordering of the nodes, and specific coordinate system.
        To meet these requirements, we can use graph neural network with embedding vector $\textbf{r}_n$, with a slightly different Aggregate function and Update function;
        we replace the place of node embeddings with relative distance between nodes such as squared distances. For more information of these functions, see Satorras, Hoogeboom, and Welling, 2021 or Bishop, 2023.
        
        Capturing such geometric symmetric features and taking it into account of the structure of the network forms the basis of a rich field of
        \textbf{geometric deep learning} (Bronstein \emph{et al.}, 2017; Bronstien \emph{et al.}, 2021).
\end{document}